%! Reference — Unit 2 — Transformations of Functions — one-page info sheet
% STANDALONE handout. Merged into no packet — see ../Makefile.
% ONE PAGE. No name row, no key: it IS the reference. Pairs with Lesson 2.2's notes: same form,
% g(x) = a f(b(x - h)) + k, same point rule (x/b + h, ay + k), same "inside runs backwards" slogan.
% Each constant gets one card in its own colour (a plum, b gold, h green, k red) with a family
% plot showing what happens AS THAT CONSTANT CHANGES — dashed slate is always the parent.
\documentclass[10pt]{article}
\usepackage{precalculus-article}
\usepackage{precalculus-boxes}

\newcommand{\ca}{{\color{plum}a}}
\newcommand{\cb}{{\color{goldacc}b}}
\newcommand{\ch}{{\color{greenacc}h}}
\newcommand{\ck}{{\color{redacc}k}}

% A family plot: axes on the window [-4,4] x [-3,5] with the body clipped to it.
\newcommand{\famplot}[1]{%
  \begin{tikzpicture}[x=0.24cm, y=0.24cm]
    \draw[linegray, very thin, step=1] (-4,-3) grid (4,5);
    \draw[->, thin] (-4.3,0) -- (4.5,0) node[right, font=\scriptsize] {$x$};
    \draw[->, thin] (0,-3.3) -- (0,5.5) node[above, font=\scriptsize] {$y$};
    \begin{scope}
      \clip (-4,-3) rectangle (4,5);
      #1
    \end{scope}
  \end{tikzpicture}}

% One constant's card: [#1 text width] #2 colour, #3 title, #4 the cases, #5 the family plot,
% #6 the caption. The full-width cards (a, b) take the default; the half-width h and k cards
% sit side by side and pass a narrower text column.
\newcommand{\constcard}[6][0.72]{%
  \begin{scenariobox}[#3]{#2}
  \begin{minipage}[c]{#1\linewidth}\small\raggedright
  #4
  \end{minipage}\hfill
  \begin{minipage}[c]{\dimexpr0.98\linewidth-#1\linewidth\relax}\centering
  #5\par\vspace{1pt}
  {\scriptsize #6\par}
  \end{minipage}
  \end{scenariobox}}

\begin{document}
\thispagestyle{empty}

\pageheader{Unit 2: Functions and Their Graphs}{Reference --- Transformations of Functions}

\vspace{-0.12in}

\begin{headlinebox}{lilac}
\centering
{\Large $g(x) \;=\; \ca\,f\big(\cb\,(x - \ch)\big) + \ck$}\par\vspace{3pt}
\small
\begin{tabular}{@{}c@{\qquad}c@{}}
\textbf{Inside} $f$: \cb\ and \ch\ act on $x$ and run \textbf{backwards}. &
\textbf{Outside} $f$: \ca\ and \ck\ act on $y$ and do \textbf{what they say}.
\end{tabular}\par\vspace{2pt}
A point $(x, y)$ on $f$ lands at
$\left(\dfrac{x}{\cb} + \ch,\ \ \ca\,y + \ck\right)$ on $g$.
\quad Dashed graph below $=$ the parent.
\end{headlinebox}

\vspace{0.03in}
\constcard{plum}{$a$ --- vertical stretch, compression, flip over the $x$-axis}{%
\textbf{Multiplies every output by $a$:} \ $(x, y) \to (x,\ a\,y)$.\par\vspace{2pt}
\begin{tabular}{@{}l@{\ \ }l@{}}
$|a| > 1$ & \textbf{vertical stretch} by $|a|$ --- taller, steeper \\
$0 < |a| < 1$ & \textbf{vertical compression} by $|a|$ --- flatter \\
$a < 0$ & also \textbf{reflects over the $x$-axis} (opens the other way) \\
$a = 1$ & no change
\end{tabular}\par\vspace{2pt}
\textbf{As $|a|$ grows} the graph climbs faster; \textbf{as $a \to 0$} it flattens onto the
$x$-axis; \textbf{crossing 0} flips it. Points on the $x$-axis ($y = 0$) never move.}%
{\famplot{%
  \draw[thick, slate, dashed, domain=-2.3:2.3, samples=40] plot (\x,{\x*\x});
  \draw[very thick, plum, domain=-1.6:1.6, samples=40] plot (\x,{2*\x*\x});
  \draw[very thick, plumlight!60, domain=-3.2:3.2, samples=40] plot (\x,{0.5*\x*\x});
  \draw[very thick, redacc, domain=-1.8:1.8, samples=40] plot (\x,{-\x*\x});
  \node[font=\tiny, plum, left] at (-1.3,3.4) {$2$};
  \node[font=\tiny, slate, right] at (1.85,3.4) {$1$};
  \node[font=\tiny, plumlight, right] at (2.6,3.4) {$\tfrac12$};
  \node[font=\tiny, redacc, right] at (1.6,-2.6) {$-1$};}}%
{$y = a\,x^2$ for $a = 2,\ 1,\ \tfrac12,\ -1$}

\constcard{goldacc}{$b$ --- horizontal compression, stretch, flip over the $y$-axis}{%
\textbf{Divides every input by $b$:} \ $(x, y) \to \left(\tfrac{x}{b},\ y\right)$.
\textbf{Backwards:} big $b$ squeezes.\par\vspace{2pt}
\begin{tabular}{@{}l@{\ \ }l@{}}
$|b| > 1$ & \textbf{horizontal compression} by $\tfrac{1}{|b|}$ --- toward the $y$-axis \\
$0 < |b| < 1$ & \textbf{horizontal stretch} by $\tfrac{1}{|b|}$ --- away from the $y$-axis \\
$b < 0$ & also \textbf{reflects over the $y$-axis} (left $\leftrightarrow$ right) \\
$b = 1$ & no change
\end{tabular}\par\vspace{2pt}
\textbf{As $|b|$ grows} the graph squeezes toward the $y$-axis; \textbf{as $b \to 0$} it
spreads out wide. Points on the $y$-axis ($x = 0$) never move.}%
{\famplot{%
  \draw[thick, slate, dashed, domain=0:4, samples=60] plot (\x,{sqrt(\x)});
  \draw[very thick, goldacc, domain=0:4, samples=60] plot (\x,{sqrt(4*\x)});
  \draw[very thick, goldacc!45, domain=0:4, samples=60] plot (\x,{sqrt(\x/4)});
  \draw[very thick, redacc, domain=-4:0, samples=60] plot (\x,{sqrt(-\x)});
  \fill[goldacc] (1,2) circle (1.5pt); \fill[slate] (4,2) circle (1.5pt);
  \node[font=\tiny, goldacc, above] at (2.6,3.2) {$4$};
  \node[font=\tiny, slate, above] at (3.4,1.85) {$1$};
  \node[font=\tiny, goldacc!80!black, below] at (3.3,0.95) {$\tfrac14$};
  \node[font=\tiny, redacc, above] at (-3.2,1.8) {$-1$};}}%
{$y = \sqrt{b\,x}$ for $b = 4,\ 1,\ \tfrac14,\ -1$}

\noindent
\begin{minipage}[t]{0.495\linewidth}\vspace{0pt}
\constcard[0.6]{greenacc}{$h$ --- horizontal shift}{%
\textbf{Adds $h$ to every input:}\\ $(x, y) \to (x + h,\ y)$. \textbf{Backwards.}\par\vspace{2pt}
\begin{tabular}{@{}l@{\ \ }l@{}}
$h > 0$ & \textbf{right} $h$: \ $f(x - 2)$ \\
$h < 0$ & \textbf{left} $|h|$: \ $f(x + 2)$ \\
$h = 0$ & no change
\end{tabular}\par\vspace{2pt}
\textbf{As $h$ increases} the graph slides right. The vertex lands where the inside is $0$.}%
{\famplot{%
  \draw[thick, slate, dashed, domain=-2.3:2.3, samples=40] plot (\x,{\x*\x});
  \draw[very thick, greenacc, domain=-0.3:4, samples=40] plot (\x,{(\x-2)*(\x-2)});
  \draw[very thick, greenacc!45, domain=-4:0.3, samples=40] plot (\x,{(\x+2)*(\x+2)});
  \fill[greenacc] (2,0) circle (1.5pt); \fill[greenacc!45] (-2,0) circle (1.5pt);
  \node[font=\tiny, greenacc, below] at (2,-0.1) {$2$};
  \node[font=\tiny, greenacc!80!black, below] at (-2,-0.1) {$-2$};}}%
{$y = (x - h)^2$,\\ $h = -2,\ 0,\ 2$}
\end{minipage}\hfill
\begin{minipage}[t]{0.495\linewidth}\vspace{0pt}
\constcard[0.6]{redacc}{$k$ --- vertical shift}{%
\textbf{Adds $k$ to every output:}\\ $(x, y) \to (x,\ y + k)$. \textbf{Does what it says.}\par\vspace{2pt}
\begin{tabular}{@{}l@{\ \ }l@{}}
$k > 0$ & \textbf{up} $k$: \ $f(x) + 3$ \\
$k < 0$ & \textbf{down} $|k|$: \ $f(x) - 3$ \\
$k = 0$ & no change
\end{tabular}\par\vspace{2pt}
\textbf{As $k$ increases} the graph slides up. The vertex, asymptote, and range move by $k$.}%
{\famplot{%
  \draw[thick, slate, dashed, domain=-2.3:2.3, samples=40] plot (\x,{\x*\x});
  \draw[very thick, redacc, domain=-1.8:1.8, samples=40] plot (\x,{\x*\x+2});
  \draw[very thick, redacc!45, domain=-2.7:2.7, samples=40] plot (\x,{\x*\x-2});
  \fill[redacc] (0,2) circle (1.5pt); \fill[redacc!45] (0,-2) circle (1.5pt);
  \node[font=\tiny, redacc, left] at (0,2) {$2$};
  \node[font=\tiny, redacc!80!black, left] at (0,-2) {$-2$};}}%
{$y = x^2 + k$,\\ $k = -2,\ 0,\ 2$}
\end{minipage}

\vspace{0.04in}
\noindent
\begin{minipage}[t]{0.58\linewidth}\vspace{0pt}
\begin{scenariobox}[All Four at Once]{plumlight}
\small\raggedright
\stepnum{1}\ \textbf{Inside:} divide $x$ by \cb, then add \ch. \quad
\stepnum{2}\ \textbf{Outside:} multiply $y$ by \ca, then add \ck.\par\vspace{3pt}
$g(x) = 3\sqrt{\tfrac12(x + 4)} - 1$: \ $\ca = 3$, $\cb = \tfrac12$, $\ch = -4$, $\ck = -1$.
Left 4, horizontal stretch by 2, vertical stretch by 3, down 1.\par\vspace{2pt}
\centering
{\renewcommand{\arraystretch}{1.15}\footnotesize
\begin{tabular}{|c|c|c|}
\hline
on $y = \sqrt{x}$ & $\tfrac{x}{b} + h$ & $a\,y + k$ \\ \hline
$(0, 0)$ & $0 \div \tfrac12 - 4 = -4$ & $3(0) - 1 = -1$ \\
$(1, 1)$ & $1 \div \tfrac12 - 4 = -2$ & $3(1) - 1 = 2$ \\
$(4, 2)$ & $4 \div \tfrac12 - 4 = 4$ & $3(2) - 1 = 5$ \\ \hline
\end{tabular}}
\end{scenariobox}
\end{minipage}\hfill
\begin{minipage}[t]{0.40\linewidth}\vspace{0pt}
\begin{scenariobox}[Watch Out]{goldacc}
\small\raggedright
\textbf{Factor out $b$ first.} \ $f(2x - 6) = f\big(2(x - 3)\big)$, so $h = 3$, not $6$.
\par\vspace{3pt}
\textbf{Minus means right.} \ $(x - 5)^2$ moves \emph{right} 5; $(x + 5)^2$ moves \emph{left}.
\par\vspace{3pt}
\textbf{Big $b$ squeezes.} \ $f(3x)$ is a compression by $\tfrac13$, not a stretch by 3.
\par\vspace{3pt}
\textbf{Negative $a$ flips up/down; negative $b$ flips left/right.}
\end{scenariobox}
\end{minipage}

\end{document}
